\documentclass[11pt]{article}
\usepackage[margin=1.1in]{geometry}
\usepackage{amsmath,amssymb,amsthm}
\usepackage{booktabs}
\usepackage[colorlinks=true,linkcolor=blue,citecolor=blue,urlcolor=blue]{hyperref}
% expansion=false: TinyTeX's default fonts are not scalable, and auto expansion
% is a fatal error there rather than a downgrade.
\usepackage[expansion=false]{microtype}

\newcommand{\R}{\mathbb{R}}
\newcommand{\E}{\mathbb{E}}
\newcommand{\Cov}{\operatorname{Cov}}
\newcommand{\tr}{\operatorname{tr}}
\newcommand{\muraw}{\theta}
\newcommand{\Sig}{\Sigma_1}

\title{Subspaces, whitening, and what a corpus can tell you\\
\large Linear algebra behind the VPop dimension reduction}
\author{}
\date{2026-08-11}

\begin{document}
\maketitle

\begin{abstract}
\noindent
These are working notes on the linear algebra used to decide that the pdac
corpus constrains roughly fifteen directions rather than 271 parameters, to
build a basis for those directions, and to carry the discarded uncertainty
forward without introducing bias. Everything reduces to three ideas: put the
problem in coordinates where both the prior and the noise are isotropic, take
an SVD there, and notice that the singular values are signal-to-noise ratios.
The rest follows.
\end{abstract}

\section{The object we are decomposing}

The model is a map from parameters to predicted data,
\[
  \tau : \R^{P} \longrightarrow \R^{K},
  \qquad P = 271 \text{ parameters}, \quad K = 138 \text{ rows}.
\]
Each row is one number a paper printed (a median, an IQR, a quantile). The
Jacobian $J = \partial\tau/\partial\mu \in \R^{K\times P}$ says how each printed
number responds to each parameter.

$J$ on its own is nearly useless, because its entries carry units. Row $i$ might
be a cell density in cells/mm$^2$ and row $i'$ a dimensionless fraction; column
$j$ might be a rate in day$^{-1}$ and column $j'$ a volume. Comparing entries of
$J$ compares apples to volts. Every step below is, at bottom, about choosing the
right units so that comparison becomes meaningful.

\subsection{Prior coordinates: making the input isotropic}

The prior on the parameter centre is Gaussian and correlated,
\[
  \mu \sim \mathcal N(\mu_0, \Sig), \qquad \Sig = L L^{\top},
\]
where $L$ is a Cholesky factor. Rather than sample $\mu$ directly, the model
samples a standard normal $\muraw$ and maps it forward:
\[
  \mu = \mu_0 + L\muraw, \qquad \muraw \sim \mathcal N(0, I_P).
\]
This is a change of variables that makes the prior isotropic: a unit step in any
direction of $\muraw$ is one prior standard deviation, and all directions are
equally likely a priori. It is the input-side whitening.

By the chain rule, the Jacobian in these coordinates is
\begin{equation}
  A \;=\; \frac{\partial \tau}{\partial \muraw}
    \;=\; \frac{\partial \tau}{\partial \mu}\,
          \frac{\partial \mu}{\partial \muraw}
    \;=\; J L .
  \label{eq:A}
\end{equation}
The columns of $A$ now all mean the same thing: ``move one prior standard
deviation along this direction, and the rows move by this much.''

\medskip
\noindent\textbf{Why not just scale each column by its marginal prior sd?}
That is what one does instinctively, and it is what we did at first. Scaling
column $j$ by $\mathrm{sd}(\mu_j)$ is the same as replacing $L$ by
$\operatorname{diag}(\mathrm{sd})$, which is correct only when $\Sig$ is
diagonal. Ours is not: 66 parameters sit in a copula block from jointly-fitted
submodel posteriors. Using $L$ keeps those correlations; using the marginal sds
silently discards them.

\subsection{Data coordinates: making the output isotropic}

The likelihood is Gaussian on the residual $r = y - \tau$ with block covariance
$V$,
\[
  -2\log p \;=\; r^{\top} V^{-1} r + \text{const}.
\]
So the natural geometry on the output side is not Euclidean; distance is
measured by $V^{-1}$. Factor $V = L_V L_V^{\top}$ and define $\tilde r = L_V^{-1} r$.
Then
\[
  r^{\top} V^{-1} r \;=\; \tilde r^{\top} \tilde r,
\]
and in the tilde coordinates the problem is ordinary least squares. This is the
output-side whitening. Applying it to the Jacobian gives the object we actually
decompose:
\begin{equation}
  \boxed{\;S \;=\; L_V^{-1} A \;=\; L_V^{-1} J L \;\in\; \R^{K\times P}.\;}
  \label{eq:S}
\end{equation}

$S$ is dimensionless in both directions. Entry $S_{ij}$ answers: \emph{if
parameter direction $j$ moves one prior standard deviation, by how many noise
standard deviations does row $i$ move?} That is a signal-to-noise ratio, and it
is the only quantity in this problem that can be compared across rows and across
parameters at once.

\medskip
\noindent\textbf{The mistake worth avoiding.} Dividing each row of $A$ by its own
standard deviation $\sqrt{V_{ii}}$ whitens only the \emph{diagonal} of $V$. Our
$V$ has substantial off-diagonal structure (rows of one readout are correlated;
blocks run to 68 rows). Section~\ref{sec:bias} shows that this distinction is
not cosmetic: it decides whether the reduction is unbiased.

\section{The SVD, and why singular values are the answer}

Take the singular value decomposition
\begin{equation}
  S \;=\; U \Sigma V^{\top},
  \qquad U^{\top}U = I, \quad V^{\top}V = I,
  \qquad \Sigma = \operatorname{diag}(\sigma_1 \ge \sigma_2 \ge \cdots).
  \label{eq:svd}
\end{equation}
Read the three factors as follows.

\begin{itemize}
  \item $v_j \in \R^{P}$: a \emph{direction in parameter space}, in units of
        prior standard deviations. It is generally a combination of many
        parameters, not a single one.
  \item $u_j \in \R^{K}$: the \emph{pattern across the data rows} that moving
        along $v_j$ produces, in units of noise standard deviations.
  \item $\sigma_j$: how many noise sds of signal you get per prior sd of
        parameter movement. A pure signal-to-noise ratio.
\end{itemize}

\subsection{The formula that makes ``identifiability'' precise}

Take the linear-Gaussian case, which is what the local analysis approximates:
prior $\muraw \sim \mathcal N(0, I)$, data $\tilde y = S\muraw + \varepsilon$
with $\varepsilon \sim \mathcal N(0, I)$. Completing the square gives the
posterior precision
\[
  \Lambda_{\text{post}} \;=\; \underbrace{I}_{\text{prior}}
    \;+\; \underbrace{S^{\top}S}_{\text{data}} ,
\]
and $S^{\top}S = V \Sigma^2 V^{\top}$ is diagonal in the $v_j$ basis. So along
direction $v_j$,
\begin{equation}
  \boxed{\;
    \frac{\text{posterior variance}}{\text{prior variance}}
    \;=\; \frac{1}{1+\sigma_j^{2}} \; }
  \label{eq:shrink}
\end{equation}

This single formula is the whole identifiability story:

\begin{center}
\begin{tabular}{ccl}
  \toprule
  $\sigma_j$ & variance ratio & reading \\
  \midrule
  $10$   & $0.01$ & the data essentially determine this direction \\
  $3$    & $0.10$ & strongly informed \\
  $1$    & $0.50$ & the data are worth as much as the prior \\
  $0.3$  & $0.92$ & barely moved \\
  $0.1$  & $0.99$ & posterior $=$ prior; invisible \\
  \bottomrule
\end{tabular}
\end{center}

A direction with $\sigma_j \ll 1$ is one where moving a full prior standard
deviation produces less than one noise standard deviation of signal. No amount
of sampling recovers it, because there is nothing to recover: the likelihood is
flat there and the posterior is the prior. Chains disagree about such a
direction not because the sampler is bad but because there is no information to
agree on.

\medskip
\noindent\textbf{Effective number of parameters.} Summing the per-direction
information gives the standard ridge-regression quantity
\begin{equation}
  P_{\text{eff}} \;=\; \sum_j \frac{\sigma_j^{2}}{1+\sigma_j^{2}},
  \label{eq:peff}
\end{equation}
the ``effective degrees of freedom''. It counts directions weighted by how much
each was actually learned, and it needs no arbitrary cutoff. If you remember one
number from a spectrum, remember this one.

\subsection{The cruder measures we quoted, and their caveat}

In practice we reported cumulative Frobenius variance,
\[
  \frac{\sum_{j \le k}\sigma_j^{2}}{\sum_j \sigma_j^{2}},
\]
answering ``what fraction of $\|S\|_F^2$ do the first $k$ directions carry''
(90\% at $k=12$, 95\% at $k=16$), and the participation ratio
\[
  \mathrm{PR} \;=\; \frac{\big(\sum_j \sigma_j^2\big)^2}{\sum_j \sigma_j^4},
\]
a cutoff-free measure of how many terms the spectrum is spread over
($\mathrm{PR} = 1$ for a single spike, $= n$ for $n$ equal values). We measured
$\mathrm{PR} \approx 5.6$.

These are \emph{not} the same question as \eqref{eq:peff}. Frobenius weighting
counts $\sigma^2$, so it is dominated by directions you already know well;
$P_{\text{eff}}$ saturates at 1 per direction no matter how large $\sigma$ gets.
A long tail of small $\sigma_j$ contributes almost nothing to Frobenius variance
and almost nothing to $P_{\text{eff}}$, but as Section~\ref{sec:prop} shows, it
can still contribute a great deal to the \emph{predictive variance} of the rows.
That is exactly the asymmetry we ran into.

\section{Two ways to pick a reduced set}

\subsection{SVD: the best $k$-dimensional subspace}

The Eckart--Young theorem says the truncated SVD is optimal: among all rank-$k$
matrices, $S_k = \sum_{j\le k}\sigma_j u_j v_j^{\top}$ minimises
$\|S - S_k\|$ in both the Frobenius and spectral norms. So if you want the best
$k$-dimensional \emph{subspace}, the leading right singular vectors are it, and
there is nothing to argue about.

The cost is interpretability: $v_j$ is a combination. ``Direction 1 is
$-0.80\,k_{\text{CD8\_T\_pro}} + 0.23\,k_{\text{CD8\_basal\_exh}} + \cdots$'' is
a real statement about the model, but it is not a parameter you can look up.

\subsection{QR with column pivoting: picking actual parameters}

If you want genuine parameters, you are asking for \emph{subset selection}:
choose $k$ columns of $S$ whose span best approximates the column space. That
problem is NP-hard in general. Column-pivoted QR is the standard greedy
heuristic:
\[
  S \Pi \;=\; QR, \qquad \Pi \text{ a permutation},
\]
where at each step the algorithm picks the remaining column with the largest
residual norm after projecting out the columns already chosen, then
orthogonalises against them. So it greedily takes ``the most informative column
that is most orthogonal to what I already have'', which is precisely what you
want from a representative subset.

Two properties worth knowing. The pivot order is a deterministic property of $S$,
so the resulting list is reproducible rather than hand-picked, which matters if a
reviewer asks how the parameters were chosen. And it is suboptimal: our first 20
pivots reproduced 91.7\% of $\|S\|_F^2$ where the SVD's first 20 directions
capture 97.4\% by construction. You pay for interpretability.

\section{Restricting to a subspace}

Let $B \in \R^{P\times k}$ have orthonormal columns, $B^{\top}B = I_k$. Two
standard objects:
\[
  \Pi \;=\; BB^{\top} \quad (\text{projector onto } \operatorname{span} B),
  \qquad
  \Pi^{\perp} \;=\; I - BB^{\top}.
\]
Both are symmetric and idempotent ($\Pi^2 = \Pi$), which is what makes the
algebra below collapse so cleanly.

The reduction is: instead of sampling $\muraw \in \R^{P}$, sample coefficients
$c \in \R^{k}$ and set
\begin{equation}
  \muraw \;=\; Bc, \qquad c \sim \mathcal N(0, I_k).
  \label{eq:restrict}
\end{equation}

\subsection{Why $B$ must be orthonormal}

Under \eqref{eq:restrict},
\[
  \Cov(\muraw) \;=\; B\,\Cov(c)\,B^{\top} \;=\; BB^{\top} \;=\; \Pi .
\]
That is exactly the prior $\mathcal N(0, I)$ restricted to the span: unit
variance in every direction inside it, zero outside. If $B$ were not
orthonormal, $BB^{\top}$ would be some other ellipsoid, and the implied prior
would be a distribution nobody chose and no one would notice. This is why the
code checks $\|B^{\top}B - I\|_\infty$ and refuses rather than proceeding.

\subsection{Why subspaces work where fixing coordinates fails}

The instinct is to hold uninformative parameters at their prior means. Setting
$\muraw_j = 0$ gives
\[
  \mu_j \;=\; \mu_{0,j} + \sum_{m} L_{jm}\,\muraw_m ,
\]
which equals $\mu_{0,j}$ \emph{only if} $L_{jm} = 0$ for all $m \ne j$, i.e.\ if
row $j$ of $\Sig$ has no off-diagonal. Otherwise the free parameters drag the
``held'' one through the correlation and it is not held at all. On our corpus
this blocks the majority of the parameters one would want to fix.

A subspace restriction makes no claim about any individual coordinate. It says
$\muraw \in \operatorname{span}(B)$, and $L$ then maps that forward with all
correlations intact. This is why $\mu$ still moved in 270 of 271 coordinates in
a fit that sampled only 16 numbers.

\section{Propagating the discarded uncertainty}
\label{sec:prop}

Restricting to the span discards the complement, which under the prior had
distribution $\mathcal N(0, \Pi^{\perp})$. Holding it at zero rather than
integrating over it throws away real epistemic uncertainty. The linear
propagation of uncertainty (the delta method) puts it back.

For a differentiable map $y = f(x)$ with $x \sim \mathcal N(m, C)$,
\[
  y \approx f(m) + \nabla f\,(x-m)
  \quad\Longrightarrow\quad
  \Cov(y) \approx \nabla f \; C \; \nabla f^{\top}.
\]
Here $\nabla f = A$ from \eqref{eq:A} and $C = \Pi^{\perp}$, so the discarded
directions contribute
\begin{equation}
  V_{\text{trunc}}
    \;=\; A\,\Pi^{\perp}\,\Pi^{\perp\top} A^{\top}
    \;=\; A\,(I - BB^{\top})\,A^{\top},
  \label{eq:trunc}
\end{equation}
using idempotence and symmetry to collapse the two projectors into one. Add this
to $V$ and the reduced model reports intervals that admit the dropped directions
were never measured, instead of intervals belonging to a model that knows them
exactly.

Note $A$ here, not $S$: the term is added to $V$, which lives in the rows' own
units, so it must not be whitened.

\medskip
\noindent\textbf{Asserted versus dropped.} One subtlety with real consequences.
A direction dropped for economy is one nobody measured, and its prior variance
belongs in $V$. But a parameter held deliberately because it is a
\emph{conditioning input} the data are not entitled to infer (for us,
\texttt{initial\_tumour\_diameter}, which sets when each patient is observed) is
\emph{asserted}. Inflating $V$ for it would claim uncertainty about a number the
model states. Those coordinates are therefore excluded from $\Pi^{\perp}$ before
forming \eqref{eq:trunc}.

\medskip
\noindent\textbf{Magnitude.} A small share of the \emph{spectrum} can be a large
share of the \emph{predictive variance}. At $k=16$ the discarded directions were
4.6\% of $\|S\|_F^2$ but inflated the median row standard deviation by 27\%,
with a maximum of $3.3\times$. There are 255 discarded directions; each is
individually negligible and they add.

\section{Bias, and why the metric decides it}
\label{sec:bias}

Inflating $V$ widens intervals. It says nothing about whether the centre is in
the right place. Truncation is a form of \textbf{omitted-variable bias}, and it
deserves separate treatment.

Split the Jacobian into retained and discarded parts,
\[
  A_{\parallel} = A\Pi, \qquad A_{\perp} = A\Pi^{\perp}, \qquad
  A = A_{\parallel} + A_{\perp}.
\]
Suppose the truth is $r = A_{\parallel}c_{\parallel} + A_{\perp}c_{\perp} + \varepsilon$
but we fit only $c_{\parallel}$, by weighted least squares with weight matrix
$W$. Then
\[
  \hat c_{\parallel}
    = (A_{\parallel}^{\top} W A_{\parallel})^{-1} A_{\parallel}^{\top} W r,
  \qquad
  \E[\hat c_{\parallel}]
    = c_{\parallel}
    + \underbrace{(A_{\parallel}^{\top} W A_{\parallel})^{-1}
       A_{\parallel}^{\top} W A_{\perp}}_{\text{bias}}\, c_{\perp}.
\]
The retained directions absorb whatever the discarded ones would have explained.
The bias vanishes exactly when
\begin{equation}
  A_{\parallel}^{\top}\,W\,A_{\perp} \;=\; 0 .
  \label{eq:nobias}
\end{equation}

Now the payoff. The likelihood's weight is $W = V^{-1}$. Using $S = L_V^{-1}A$
and its SVD \eqref{eq:svd} with $B = V_k$ (the leading right singular vectors),
\[
  S\Pi = U_k \Sigma_k V_k^{\top}\!\cdot\!(\text{restricted}),
  \qquad
  S\Pi^{\perp} = U_{>k}\Sigma_{>k}V_{>k}^{\top}\!\cdot\!(\text{restricted}),
\]
and therefore
\[
  A_{\parallel}^{\top} V^{-1} A_{\perp}
    = (L_V^{-1}A_{\parallel})^{\top}(L_V^{-1}A_{\perp})
    = (S\Pi)^{\top}(S\Pi^{\perp})
    \;\propto\; \Sigma_k\, U_k^{\top} U_{>k}\, \Sigma_{>k}
    \;=\; 0,
\]
because the columns of $U$ are orthonormal. \textbf{The SVD taken in the
whitened metric makes the omitted directions orthogonal to the retained ones in
observation space, so the omitted-variable bias is exactly zero at first order.}

This is the reason the metric matters, and it is not a refinement. Take the SVD
in any other metric and \eqref{eq:nobias} holds only approximately. Under
row-standard-deviation scaling it holds to the extent $V$ is diagonal, which
ours is not. And note that coordinate selection (QR, or holding parameters)
enjoys no such protection at all: retained and dropped \emph{columns} are not
generally orthogonal in any metric, so subset selection biases at first order
while the subspace does not. That is an argument for the SVD basis independent
of everything else.

\subsection{What survives}

The cancellation is first order. The model is nonlinear, so
$\E[\tau(\mu)] \ne \tau(\E[\mu])$, and a second-order expansion gives
\[
  \E[\tau(\mu)] \;\approx\; \tau(\bar\mu) + \tfrac12 \tr\!\big(H\,C\big),
\]
with $H$ the Hessian of $\tau$. Removing the complement changes $C$ from $I$ to
$\Pi$, shifting the predicted mean by $-\tfrac12\tr(H\,\Pi^{\perp})$. Equation
\eqref{eq:trunc} does not capture this. It is small when the discarded variance
is small, but it is an approximation and should be stated as one rather than
claimed as an identity.

\section{Summary of the pipeline}

\begin{enumerate}
  \item Compute $J = \partial\tau/\partial\mu$ by automatic differentiation.
  \item Right-multiply by $L$ (prior Cholesky) to get $A = JL$: input whitening,
        so a unit step is one prior sd and correlations are respected.
  \item Left-multiply by $L_V^{-1}$ (noise Cholesky, block by block) to get
        $S = L_V^{-1}A$: output whitening, so a unit response is one noise sd.
  \item SVD $S = U\Sigma V^{\top}$. The $\sigma_j$ are signal-to-noise ratios;
        \eqref{eq:shrink} converts them into how much each direction is learned.
  \item Keep $B = V_k$ for a $k$ chosen from the spectrum. Sample
        $c\in\R^k$ with $\muraw = Bc$.
  \item Add $A(I-BB^{\top})A^{\top}$ to $V$ so the discarded directions widen the
        error bars instead of vanishing, excluding any coordinate that is
        asserted rather than dropped.
  \item Bias is zero at first order by construction, \emph{because} step 3 used
        the likelihood's own metric.
\end{enumerate}

\section{Where this sits in the literature}

None of this is new, and the names are worth knowing for the write-up.

\begin{itemize}
  \item \textbf{Likelihood-informed subspace (LIS).} Cui, Martin, Marzouk,
        Solonen and Spantini (2014), \emph{Likelihood-informed dimension
        reduction for nonlinear inverse problems}, Inverse Problems 30(11).
        This is essentially the method above: whiten by the prior, whiten by the
        noise, take the dominant eigendirections of the resulting operator.
  \item \textbf{Optimal low-rank posterior approximation.} Spantini et al.\
        (2015), SIAM J.\ Sci.\ Comput.\ 37(6), proves in what sense the leading
        directions of \eqref{eq:svd} are the optimal reduction, which is the
        rigorous version of the hand-waving in Section~3.
  \item \textbf{Active subspaces.} Constantine (2015). Closely related but
        prior-driven rather than likelihood-driven: it averages
        $\nabla f \nabla f^{\top}$ over an input distribution. The distinction
        matters, since a direction can be influential and still unidentified.
  \item \textbf{Ridge regression and effective degrees of freedom.} Hastie,
        Tibshirani and Friedman, \emph{ESL}, §3.4. Equation \eqref{eq:peff} is
        the same quantity; a Gaussian prior is Tikhonov regularisation, and
        \eqref{eq:shrink} is the shrinkage factor.
  \item \textbf{Column-pivoted QR and subset selection.} Golub and Van Loan,
        \emph{Matrix Computations}, §5.4.2.
  \item \textbf{Sloppy models.} Gutenkunst et al.\ (2007), PLoS Comput Biol
        3(10). The systems-biology observation that these spectra are generically
        near-exponential, which is exactly what we see, and why the effective
        rank is small for reasons that are structural rather than a defect of
        this corpus.
\end{itemize}

\end{document}
